\section{Related work}
\label{sec:related-work}

The work touches three literatures: the statistics that validate probability of default models and the instructions that fix them, the models and numerical methods of portfolio credit capital, and the interest rate risk and model-risk toolkits that a validator can install. Each paragraph below says what a group of methods assumes, what it gives and where it stops, and the last paragraph of the section says where the engine stands.

\subsection{Validation statistics and the reporting instructions}

The binomial test of a grade's default count against its probability of default is the oldest back-test and assumes independent defaults; the Basel Committee's study of validation shows that under a common factor its true size can be much larger than the nominal one, and it gives a critical value corrected for the asset correlation~\citep{bcbswp14}. The Jeffreys test replaces the binomial tail by the posterior probability, under the Beta prior with both parameters one half, that the true probability is below the applied one; its interval has good coverage for small samples~\citep{brown2001} and the ECB instructions adopt it as the reported back-test of every grade and of the portfolio~\citep{ecb2019}. For discrimination the area under the ROC curve has a variance that follows from the structural components of the Mann--Whitney statistic~\citep{hanley1982,delong1988}; the instructions use the variance of those components to test whether the current area is below the area at development. For stability the population stability index has a chi-square null with a sampling distribution worked out by Yurdakul and Naranjo~\citep{yurdakul2020}, which the engine reproduces beside the ECB's own concentration test; and the tests of a rating migration matrix, whether it is homogeneous in time or of first order, are the Anderson and Goodman chi-square tests~\citep{anderson1957}. For market risk the coverage tests of Kupiec are the model of a test whose size and power were studied from the start~\citep{kupiec1995}. What this literature does not provide is an executable catalogue: each paper gives a statistic, and the choices that make a test the one a supervisor means (the weighting of the portfolio probability, the tail, the missing cells) live in an instruction document.

\subsection{Portfolio credit models and their numerical methods}

The structural model of default as the crossing of a threshold by the value of the firm is Merton's~\citep{merton1974}; with one common factor its large-portfolio limit has a closed-form loss distribution that underlies the Basel risk-weight functions~\citep{gordy2003}, and the Gaussian copula is the standard way to couple defaults in a finite portfolio~\citep{li2000}. A comparative study of the commercial models, a structural one and an actuarial one, shows that they reduce to common mathematics once the inputs are matched~\citep{gordy2000}; the heavy-tailed alternatives, such as the $t$ copula, are analysed in~\citep{frey2003}. The actuarial model, CreditRisk+, gives the loss distribution by a recursion on a grid of loss units, the compound-distribution recursion of Panjer~\citep{panjer1981}, and its numerically stable form is due to Haaf, Reiss and Schoenmakers~\citep{haaf2004}. The finite-portfolio correction to the one-factor capital, the granularity adjustment, has a regulatory form with a printed test vector~\citep{gordylutkebohmert2013}, and the capital of a name or a sector follows from the Euler principle~\citep{tasche2008}, estimated for the value at risk with kernel methods~\citep{tasche2009} and for the one-factor model in closed form~\citep{emmer2005}. Rare tail events are sampled by importance sampling with an exponential twist~\citep{glasserman2005}. The calibration of the default probability curve to default data is treated in~\citep{tasche2013} and, for portfolios with few defaults, the estimation of the default probability in~\citep{pluto2005}; the asset correlation is not calibrated by any of the works above and is estimated here by the engine and checked against a simulator with a known truth. These works give the models and, in the best cases, a number to reproduce; none gives the code, and an engine built from them is checked only as far as its builder chooses to check it.

\subsection{Toolkits and the place of the proposed engine}

The open toolkits that a validator can install cover other ground. PiML is a toolbox for interpretable machine learning model development and diagnostics, with banking examples~\citep{sudjianto2023}; Shapley attributions are the common explanation method~\citep{lundberg2017}; conformal prediction gives distribution-free intervals~\citep{angelopoulos2023}; and the Hull and White model is the standard interest-rate generator for exposure simulation~\citep{hull1990}. The survey of Section~\ref{sec:results} reads the metadata of \pypiFound{} such packages and scans the sources of \distScanned{} installed distributions for the vocabulary of the catalogue, and finds a match in words only: the Jeffreys interval inside a statistics library, the Herfindahl index inside a binning library, and nothing for the ECB migration tests, CreditRisk+, the granularity adjustment or the standardised economic value of equity.

The engine studied here stands as an implementation of those statistics and models in one package, with the registry and the result record as the interface, and the contribution of this paper is not a new statistic: the arithmetic of every test is in the cited papers. It is the set of measurements that say how far each implemented piece can be trusted, and the protocol that produced them. A related question, whether a loss function or an elicitable score is the right way to rank value-at-risk models~\citep{fissler2016}, is outside the paper's scope and is part of the same catalogue only as a registered test.
